\chapter{Quy ước benchmark và nguồn artifact}
\label{ap:benchmark}
\markboth{PHỤ LỤC \Alph{chapter}. QUY ƯỚC BENCHMARK VÀ NGUỒN ARTIFACT}{}

Prompt red-team dùng định dạng ID \texttt{RT-<NHOM>-<NNN>}, với các nhóm
\texttt{BENIGN}, \texttt{INJ-DIRECT}, \texttt{INJ-ROLE},
\texttt{INJ-HIER}, \texttt{INJ-JAILBREAK}, \texttt{INJ-EXTRACT},
\texttt{INJ-RAGCTX} và \texttt{INJ-TOOL}. Nguồn benchmark đóng băng là
\texttt{redteam/prompts.jsonl}. Nhãn mong đợi không bị thay đổi trong Phase 7.1.

Các artifact kết quả gồm \texttt{latest-evaluation.json/.md},
\texttt{failure-triage.md}, \texttt{baseline-vs-guarded.json/.md} và
\texttt{reports/evidence/evidence-index.md}.

\chapter{Cấu hình và lệnh tái lập}
\label{ap:reproduction}
\markboth{PHỤ LỤC \Alph{chapter}. CẤU HÌNH VÀ LỆNH TÁI LẬP}{}

Hệ thống chạy trong virtual environment Python với dependency trong
\texttt{requirements.txt}. Provider mặc định là \texttt{mock}, model name là
\texttt{mock-rag-guard-v1}; không cần \texttt{.env} hoặc API key.

\begin{verbatim}
python -m venv .venv
.\.venv\Scripts\Activate.ps1
python -m pip install -r requirements.txt
python -m pytest -q
python scripts\run_evaluation.py
python scripts\run_evaluation.py --comparison
\end{verbatim}

Kết quả được xác minh khi đóng gói báo cáo: 82 test pass, guarded 40/40 và
baseline 5/40. Một warning deprecation từ Starlette có thể xuất hiện; dự án
không cài \texttt{httpx2}.

\chapter{Tóm tắt kết quả theo decision}
\label{ap:decision-results}
\markboth{PHỤ LỤC \Alph{chapter}. TÓM TẮT KẾT QUẢ THEO DECISION}{}

\begin{table}[H]
\centering
\caption{Phân bố expected decision của 40 prompt}
\begin{tabular}{lr}
\toprule
Expected decision & Số lượng \\
\midrule
Allow & 5 \\
Block & 34 \\
Sanitize & 1 \\
Log only & 0 \\
Human review & 0 \\
\bottomrule
\end{tabular}
\end{table}

Sau calibration, guarded mode khớp expected decision cho toàn bộ 40 prompt.
Always-allow baseline chỉ khớp 5 prompt benign. Chi tiết từng case, matched rule,
reason và risk score nằm trong artifact JSON.

\chapter{Nguồn tham khảo cần review toàn văn}
\label{ap:pending-review}
\markboth{PHỤ LỤC \Alph{chapter}. NGUỒN THAM KHẢO CẦN REVIEW TOÀN VĂN}{}

Các nguồn sau đã được xác minh thông tin thư mục nhưng trạng thái đọc toàn văn
cần được nhóm xác nhận trước bản nộp cuối:

\begin{itemize}
    \item \cite{zou2024poisonedrag} -- PoisonedRAG;
    \item \cite{wang2026pidpattack} -- PIDP-Attack;
    \item \cite{gulyamov2026review} -- bài tổng quan prompt injection.
\end{itemize}

Không sử dụng số liệu chi tiết từ các nguồn này nếu chưa hoàn tất review.

\chapter{Danh sách hình cần hoàn thiện thủ công}
\label{ap:pending-figures}
\markboth{PHỤ LỤC \Alph{chapter}. DANH SÁCH HÌNH CẦN BỔ SUNG}{}

\begin{table}[H]
\centering
\small
\caption{Đặc tả ba hình evidence cần bổ sung trước bản nộp}
\label{tab:pending-figure-specification}
\begin{tabularx}{\textwidth}{p{3.6cm}p{3.2cm}X}
\toprule
Tệp trong \texttt{figure/} & Nhãn & Nguồn và nội dung bắt buộc \\
\midrule
\texttt{architecture-overview.png} & \texttt{fig:architecture-overview} &
\texttt{docs/diagrams/architecture.md}; phân biệt pipeline đã hiện thực và các
thành phần nằm ngoài phạm vi proof-of-concept. \\
\texttt{evaluation-result-40-cases.png} & \texttt{fig:evaluation-result} &
Đầu ra \texttt{python scripts/run\_evaluation.py}; hiển thị 40/40, 0 false
positive và 0 false negative trên controlled synthetic benchmark. \\
\texttt{baseline-vs-guarded-comparison.png} & \texttt{fig:baseline-vs-guarded} &
\texttt{reports/evaluation/baseline-vs-guarded.md}; hiển thị cả baseline và
guarded, đồng thời không diễn giải baseline như chất lượng LLM. \\
\bottomrule
\end{tabularx}
\end{table}

Caption cuối cùng đã được khai báo tại vị trí chèn hình trong Chương 2 và
Chương 4. Khi tệp không hiện diện, template hiển thị khung thông báo thiếu tệp
hình thay vì làm hỏng quá trình biên dịch. Nhóm phải xuất hình từ chính
repository/artifact và không
sử dụng ảnh minh họa giả để thay evidence.
